%% toc: P3 33 33.5,33.6,33.7,33.8
\section{Proofs for the framing and saturation results of Morley's \emph{Categoricity in Power}}
\label{sec:annex-p03-w03-p03-034}

The proofs below work in Morley's normalized setting
(\cref{def:p03-morley-standing}). Two facts are used repeatedly. First,
because \(T\) is Morleyized, every formula with parameters from a system is
equivalent modulo \(T\) to an atomic relation, so a monomorphism preserves
and reflects every formula, an inclusion of models of \(T\) is elementary,
and unions of chains of models are models
(\cref{thm:p03-morley-normalization}). Second, every point of \(S(A)\) is
realized in some model containing \(A\) (\cref{thm:p03-morley-types}); an
element realizes exactly one point of \(S(A)\).

\emph{Transport criterion.} Let \(f:D\to E\) be an isomorphism between
systems of \(\mathcal N(T)\). Let \(a\notin D\) lie in a model containing
\(D\), realizing \(p\in S(D)\), and let \(b\notin E\) lie in a model
containing \(E\). Then \(f\cup\{(a,b)\}\) is a monomorphism of the induced
system on \(|D|\cup\{a\}\) into the induced system on \(|E|\cup\{b\}\) if and
only if \(b\) realizes the transport \((f^{-1})^*(p)\). Indeed, a relation
instance among elements of \(|D|\cup\{a\}\) that does not involve \(a\) is
preserved and reflected because \(f\) is an isomorphism. One that involves
\(a\) becomes, on replacing \(a\) by \(v_0\), a formula
\(\psi(v_0,\bar d)\in F(D)\), and its image is \(\psi(v_0,f(\bar d))\). So
the extended map preserves and reflects all relations exactly when, for
every \(\psi(v_0,\bar d)\in F(D)\), \(\psi\in p\) holds if and only if \(b\)
satisfies \(\psi(v_0,f(\bar d))\). The transport consists of the formulas
\(\psi(v_0,f(\bar d))\) with \(\psi(v_0,\bar d)\in p\); since the set of
formulas satisfied by \(b\) and the transport are both points of \(S(E)\),
the one-way inclusion is already equality. Injectivity is included, since
equality is one of the relations. The same argument with \(f\) the identity
gives the realization statement of \cref{thm:p03-morley-types}. The
inverse of a monomorphism onto a system is again a monomorphism, and the
union of an increasing chain of monomorphisms is a monomorphism, because
each relation instance involves finitely many elements.

\begin{annexitem}{annex:p03-w03-p03-034-dlo-powers}{thm:p03-w03-p03-034-dlo-powers}{Two dense orders of each uncountable power}
\begin{proof}
The countable case is \cref{thm:dlo-countable-categorical}.

Let \(\kappa\) be uncountable. In \(L_\kappa\), put
\((\alpha,q)<(\beta,r)\) when \(\alpha<\beta\), or \(\alpha=\beta\) and
\(q<r\); this is a linear order. It is dense: if
\((\alpha,q)<(\beta,r)\) with \(\alpha=\beta\), the point
\((\alpha,(q+r)/2)\) lies between them, and if \(\alpha<\beta\), the point
\((\alpha,q+1)\) does. It has no endpoints, since
\((\alpha,q-1)<(\alpha,q)<(\alpha,q+1)\). Its power is
\(\kappa\cdot\aleph_0=\kappa\). Density and the absence of endpoints are
preserved by reversing the order, so \(L_\kappa^{\mathrm{rev}}\) is also a
model of \(\mathsf{DLO}\) of power \(\kappa\).

Since \(\omega_1\leq\kappa\), the elements \((\xi,0)\) for \(\xi<\omega_1\)
form a strictly increasing sequence of order type \(\omega_1\) in
\(L_\kappa\).

Suppose \((x_\xi)_{\xi<\omega_1}\) were strictly increasing in
\(L_\kappa^{\mathrm{rev}}\), that is, strictly decreasing in \(L_\kappa\),
with \(x_\xi=(\alpha_\xi,q_\xi)\). If \(\xi<\eta\) then
\(x_\eta<x_\xi\) in \(L_\kappa\), so \(\alpha_\eta\leq\alpha_\xi\). Let
\(\alpha^*\) be the least of the ordinals \(\alpha_\xi\), attained at
\(\xi_0\). For \(\xi\geq\xi_0\) we have
\(\alpha^*\leq\alpha_\xi\leq\alpha_{\xi_0}=\alpha^*\), so
\(\alpha_\xi=\alpha^*\). Hence for \(\xi_0\leq\xi<\eta<\omega_1\) the
rationals satisfy \(q_\eta<q_\xi\). The map \(\xi\mapsto q_\xi\) on the
uncountable set \(\{\xi:\xi_0\leq\xi<\omega_1\}\) is then injective into
\(\mathbb Q\), which is impossible.

An isomorphism \(L_\kappa\to L_\kappa^{\mathrm{rev}}\) would carry the
increasing \(\omega_1\)-sequence of \(L_\kappa\) to one in
\(L_\kappa^{\mathrm{rev}}\). So the two models of power \(\kappa\) are not
isomorphic, and \(\mathsf{DLO}\) is not categorical in \(\kappa\).
\end{proof}
\end{annexitem}

\begin{annexitem}{annex:p03-w03-p03-034-basis-power}{thm:p03-w03-p03-034-basis-power}{Counting the canonical basis}
\begin{proof}
The map \(\varphi\mapsto U_\varphi\) is defined on equivalence classes,
since equivalent formulas lie in the same points. It is injective: if
\(\varphi\) and \(\psi\) are not equivalent modulo \(T(A)\), one of
\(\varphi\wedge\neg\psi\) and \(\psi\wedge\neg\varphi\) is consistent with
\(T(A)\), hence is a nonzero element of \(F(A)\) and lies in some
ultrafilter \(p\in S(A)\), which belongs to exactly one of \(U_\varphi\)
and \(U_\psi\). It is onto the canonical basis by definition.

Upper bound. The language of \(T(A)\) consists of the countably many
symbols of \(T\) together with one constant for each element of \(A\).
Formulas are finite strings over this alphabet, so there are at most
\(\kappa(A)+\aleph_0\) of them, and \(F(A)\) is a quotient of a subset of
them.

Lower bound. For distinct \(a,a'\in A\), the diagram of \(A\) contains
\(\neg(a=a')\), because equality is retained. In every model of \(T(A)\)
the interpretation of \(a\) satisfies \(v_0=a\) and not \(v_0=a'\), so the
two formulas are inequivalent. Thus \(a\mapsto[v_0=a]\) is injective and
\(\kappa(A)\leq\kappa(F(A))\).

Together, \(\kappa(A)+\aleph_0\leq\kappa(F(A))+\aleph_0\leq
\kappa(A)+\aleph_0\), and by the bijection the canonical basis has the
power of \(F(A)\).

The example. Let \(\Sigma\) be the theory, in the language with equality
only, asserting for each \(n\) that there are at least \(n\) elements. Any
two countably infinite models are isomorphic and \(\Sigma\) has no finite
model, so \(\Sigma\) is complete by the completeness argument recorded in
\cref{sec:p03-morley-normalization}. In a model \(M\) every permutation is
an automorphism, also of the definitional expansion, and any two elements
are exchanged by a transposition. Hence a formula \(\varphi(v_0)\) without
parameters holds of all elements of \(M\) or of none, and by completeness
\(T\) proves \(\forall v_0\,\varphi\) or \(\forall v_0\,\neg\varphi\).
So \(F(\varnothing)\) has the two elements \([v_0=v_0]\) and
\([\neg v_0=v_0]\), and the canonical basis is
\(\{S(\varnothing),\varnothing\}\).
\end{proof}
\end{annexitem}

\begin{annexitem}{annex:p03-w03-p03-034-abstract-derivative}{thm:p03-w03-p03-034-abstract-derivative}{Closedness and exact images}
\begin{proof}
\emph{A topological proposition} \cite[p.~520]{Morley1965}. Let \(G\) be
compact, \(H\) Hausdorff, \(h:G\to H\) continuous and onto, and \(r\in H\)
not isolated in \(H\). Then \(h^{-1}(r)\) contains a point that is not
isolated in \(G\). Otherwise \(h^{-1}(r)\) is a union of open singletons,
hence open, so \(G\setminus h^{-1}(r)\) is closed in \(G\) and compact. Its
image is \(H\setminus\{r\}\), because \(h\) is onto; this image is compact,
hence closed in the Hausdorff space \(H\), and \(r\) would be isolated.

\emph{Clause 1.} By the recursion,
\(X^{\alpha+1}=X\setminus\bigcup_{\beta\leq\alpha}\operatorname{Tr}^\beta(X)
=X^\alpha\setminus\operatorname{Tr}^\alpha(X)\), and for a limit
\(\delta\), a point lies in \(X^\delta\) exactly when it lies in no
\(\operatorname{Tr}^\beta(X)\) with \(\beta<\delta\), that is, in every
\(X^{\beta+1}\) with \(\beta<\delta\); so
\(X^\delta=\bigcap_{\beta<\delta}X^\beta\), the stages being decreasing.
Applying the defining condition to the identity of \(X\), which belongs to
\(\mathcal C\), shows that every point of \(\operatorname{Tr}^\alpha(X)\)
is isolated in \(X^\alpha\). So \(\operatorname{Tr}^\alpha(X)\) is open in
\(X^\alpha\) and \(X^{\alpha+1}\) is closed in \(X^\alpha\). By induction,
with \(X^0=X\) and intersections at limits, every \(X^\alpha\) is closed in
\(X\), hence compact.

\emph{Clause 2, inclusion \(\phi(Y^\alpha)\subseteq X^\alpha\).} Let
\(q\in Y^\alpha\) and \(\beta<\alpha\). Then \(q\in Y^\beta\) and
\(q\notin\operatorname{Tr}^\beta(Y)\), so there are a map
\(\psi:V\to Y\) of \(\mathcal C\) and a point
\(r\in\psi^{-1}(q)\cap V^\beta\) that is not isolated in \(V^\beta\).
The composite \(\phi\circ\psi:V\to X\) belongs to \(\mathcal C\) and
\(r\in(\phi\circ\psi)^{-1}(\phi(q))\cap V^\beta\), so
\(\phi(q)\notin\operatorname{Tr}^\beta(X)\). As \(\beta<\alpha\) was
arbitrary, \(\phi(q)\in X^\alpha\). This inclusion holds for every ordinal.

\emph{Clause 2, inclusion \(X^\alpha\subseteq\phi(Y^\alpha)\).} Assume
clause~2 for all ordinals below \(\alpha\). Let \(p\in X^\alpha\) and
suppose \(\phi^{-1}(p)\cap Y^\alpha=\varnothing\). Put
\(Z_\beta=\phi^{-1}(p)\cap Y^\beta\) for \(\beta\leq\alpha\). Since
\(\{p\}\) is closed and \(\phi\) is continuous and onto, \(Z_0\) is closed
and nonempty; the \(Z_\beta\) are closed and decrease, and
\(Z_\alpha=\varnothing\). Let \(\beta^*\) be least with
\(Z_{\beta^*}=\varnothing\). Then \(\beta^*\neq0\), and \(\beta^*\) is not
a limit, since a decreasing family of nonempty closed subsets of the
compact space \(Y\) has nonempty intersection. So \(\beta^*=\gamma+1\)
with \(\gamma<\alpha\), \(Z_\gamma\neq\varnothing\), and by clause~1
\(Z_\gamma\subseteq\operatorname{Tr}^\gamma(Y)\). Since
\(p\in X^\alpha\subseteq X^{\gamma+1}\), we have \(p\in X^\gamma\) and
\(p\notin\operatorname{Tr}^\gamma(X)\), so there are a map
\(\psi:V\to X\) of \(\mathcal C\) and \(r\in\psi^{-1}(p)\cap V^\gamma\)
not isolated in \(V^\gamma\). By amalgamation there are \(W\),
\(\phi':W\to Y\) and \(\psi':W\to V\) in \(\mathcal C\) with
\(\phi\circ\phi'=\psi\circ\psi'\). By the induction hypothesis
\(\psi'(W^\gamma)=V^\gamma\). The restriction of \(\psi'\) to the compact
space \(W^\gamma\) is a continuous map onto the Hausdorff space
\(V^\gamma\), so by the proposition there is
\(s\in W^\gamma\cap\psi'^{-1}(r)\) that is not isolated in \(W^\gamma\).
By the first inclusion, \(\phi'(s)\in Y^\gamma\), and
\(\phi(\phi'(s))=\psi(\psi'(s))=\psi(r)=p\). Hence
\(\phi'(s)\in Z_\gamma\subseteq\operatorname{Tr}^\gamma(Y)\). Applying the
definition of \(\operatorname{Tr}^\gamma(Y)\) to the map \(\phi'\), every
point of \(\phi'^{-1}(\phi'(s))\cap W^\gamma\) is isolated in
\(W^\gamma\); but \(s\) is such a point and is not isolated. This
contradiction proves \(X^\alpha\subseteq\phi(Y^\alpha)\), and clause~2
follows by induction on \(\alpha\).

\emph{Clause 3.} Let \(p\in X^\alpha\). If
\(p\in\operatorname{Tr}^\alpha(X)\), then \(p\notin X^{\alpha+1}\) by
clause~1; a point \(q\in\phi^{-1}(p)\cap Y^{\alpha+1}\) would give
\(p=\phi(q)\in X^{\alpha+1}\) by clause~2 at \(\alpha+1\). So
\(\phi^{-1}(p)\cap Y^\alpha\subseteq Y^\alpha\setminus Y^{\alpha+1}
=\operatorname{Tr}^\alpha(Y)\). If \(p\notin\operatorname{Tr}^\alpha(X)\),
then \(p\in X^{\alpha+1}=\phi(Y^{\alpha+1})\), so some
\(q\in\phi^{-1}(p)\cap Y^{\alpha+1}\) lies in \(Y^\alpha\) but not in
\(\operatorname{Tr}^\alpha(Y)\).
\end{proof}
\end{annexitem}

\begin{annexitem}{annex:p03-w03-p03-034-cantor-bendixson}{thm:p03-w03-p03-034-cantor-bendixson}{One space, one map}
\begin{proof}
The only map is the identity, which is closed under composition; for
amalgamation take \(W=X\) and both maps the identity. For
\(p\in X^\alpha\), the defining condition of
\(\operatorname{Tr}^\alpha(X)\) concerns only the identity, for which
\(\mathrm{id}^{-1}(p)\cap X^\alpha=\{p\}\); so
\(p\in\operatorname{Tr}^\alpha(X)\) exactly when \(p\) is isolated in
\(X^\alpha\). By clause~1 of \cref{thm:p03-w03-p03-034-abstract-derivative},
\(X^{\alpha+1}\) is the set of non-isolated points of \(X^\alpha\) and
\(X^\delta=\bigcap_{\beta<\delta}X^\beta\) at limits, while \(X^0=X\).
These are the clauses defining \(X^{(\alpha)}\), so
\(X^\alpha=X^{(\alpha)}\) for every \(\alpha\) by induction.
\end{proof}
\end{annexitem}

\begin{annexitem}{annex:p03-w03-p03-034-universal-homogeneous}{thm:p03-w03-p03-034-universal-homogeneous}{Back and forth}
\begin{proof}
Let \(\kappa=\kappa(A)\), an infinite cardinal.

\emph{Saturated implies homogeneous.} Let \(B_1,B_2\subseteq A\) with
\(\kappa(B_i)<\kappa\) and let \(f:B_1\to B_2\) be an isomorphism.
Enumerate \(|A|=\{a_\xi:\xi<\kappa\}\). We build isomorphisms
\(f_\xi:D_\xi\to E_\xi\) between subsystems of \(A\), increasing in
\(\xi\le\kappa\), with \(f_0=f\), unions at limits, and
\(\kappa(D_\xi)\leq\kappa(B_1)+2\cdot\kappa(\xi)\), which is \(<\kappa\)
for \(\xi<\kappa\) because both summands are below the infinite cardinal
\(\kappa\). At a limit the union of the chain is an isomorphism by the
remark after the transport criterion. At stage \(\xi+1\), first the
forth step: if \(a_\xi\notin D_\xi\), let \(p\in S(D_\xi)\) be the point
it realizes and \(p'\) its transport along \(f_\xi\). Since
\(\kappa(E_\xi)<\kappa\), saturation gives \(b\in A\) realizing \(p'\).
As \(p\) contains \(\neg(v_0=d)\) for every \(d\in D_\xi\), the transport
contains \(\neg(v_0=e)\) for every \(e\in E_\xi\), so \(b\notin E_\xi\).
By the transport criterion \(f_\xi\cup\{(a_\xi,b)\}\) is a monomorphism
onto the induced system on \(|E_\xi|\cup\{b\}\), hence an isomorphism of
subsystems of \(A\). Then the back step: if \(a_\xi\) is not yet in the
range, apply the forth step to the inverse isomorphism and the element
\(a_\xi\), and invert the result. The union \(f_\kappa\) has domain and
range containing every \(a_\xi\), so it is an automorphism of \(A\)
extending \(f\).

\emph{Saturated implies universal.} Let \(C\in\mathcal N(T)\) with
\(\kappa(C)\leq\kappa\); \(C\) is a subsystem of some model \(M\). If \(C\)
is empty there is nothing to prove. Otherwise enumerate
\(|C|=\{c_\xi:\xi<\mu\}\), where \(\mu\) is the initial ordinal of
\(\kappa(C)\), and let \(C_\xi\) be the induced system on
\(\{c_\eta:\eta<\xi\}\); then \(\kappa(C_\xi)=\kappa(\xi)<\kappa(C)\leq
\kappa\). Build increasing monomorphisms \(g_\xi:C_\xi\to A\), with
\(g_0\) empty and unions at limits. At stage \(\xi+1\), \(g_\xi\) is an
isomorphism onto the subsystem \(E=g_\xi(C_\xi)\) of \(A\), and
\(\kappa(E)<\kappa\). Let \(p\in S(C_\xi)\) be the point realized by
\(c_\xi\) in \(M\); saturation gives \(b\in A\) realizing its transport,
and \(b\notin E\) as before. By the transport criterion
\(g_{\xi+1}=g_\xi\cup\{(c_\xi,b)\}\) is a monomorphism. The union is a
monomorphism of \(C\) into \(A\).

\emph{Universal and homogeneous implies saturated.} Let \(D\subseteq A\)
with \(\kappa(D)<\kappa\) and \(p\in S(D)\). If \(p\) contains
\(v_0=d\) for some \(d\in D\), then \(d\) realizes \(p\) in \(A\).
Otherwise let \(b\) realize \(p\) in a model containing \(D\); then
\(b\notin D\). The induced system \(C\) on \(|D|\cup\{b\}\) belongs to
\(\mathcal N(T)\) and \(\kappa(C)\leq\kappa(D)+1\leq\kappa\): if
\(\kappa(D)\) is infinite this is \(\kappa(D)\), and otherwise it is
finite. By universality there is a monomorphism \(g:C\to A\). Its
restriction to \(D\) is an isomorphism of \(D\) onto the subsystem
\(g(D)\) of \(A\), and \(\kappa(g(D))=\kappa(D)<\kappa\). By homogeneity
the inverse isomorphism \(g(D)\to D\) extends to an automorphism \(h\) of
\(A\). Then \(h\circ g\) is a monomorphism of \(C\) into \(A\) that is the
identity on \(D\) and sends \(b\) to \(h(g(b))\), which is not in \(D\) by
injectivity. By the transport
criterion with \(f\) the identity, \(h(g(b))\) realizes \(p\). So \(A\) is
saturated.
\end{proof}
\end{annexitem}

\begin{annexitem}{annex:p03-w03-p03-034-limit-power}{thm:p03-w03-p03-034-limit-power}{Categoricity in an increasing sequence of powers}
\begin{proof}
Because the sequence is strictly increasing and \(\gamma\) is a nonzero limit
ordinal, \(\kappa_1>\kappa_0\geq\aleph_0\), so \(\Sigma\) is categorical
in the uncountable cardinal \(\kappa_1\). By Morley's theorem
(\cref{thm:p03-morley-original-categoricity}) \(\Sigma\) is categorical in
every uncountable cardinal. The supremum \(\lambda\) is at least
\(\kappa_1\), hence uncountable, so \(\Sigma\) is categorical in
\(\lambda\). Categoricity of \(\Sigma\) and of \(T=\Sigma^*\) are
equivalent by \cref{thm:p03-morley-normalization}, so the statement may be
read for either. No hypothesis on cardinal exponentiation was used.
\end{proof}
\end{annexitem}

\begin{annexitem}{annex:p03-w03-p03-034-mv-existence}{thm:p03-w03-p03-034-mv-existence}{An \(\omega_1\)-chain when \(\kappa^{\aleph_0}=\kappa\)}
\begin{proof}
Note that \(2^{\aleph_0}\leq\kappa^{\aleph_0}=\kappa\).

\emph{One step.} Let \(B\) be a model of \(T\) of power \(\kappa\). We find a
model \(B'\supseteq B\) of power \(\kappa\) in which every point of
\(S(A)\) is realized, for every countable \(A\subseteq B\). There are at
most \(\kappa^{\aleph_0}=\kappa\) countable subsets of \(|B|\). For each,
\(F(A)\) is countable by \cref{thm:p03-w03-p03-034-basis-power}, so
\(S(A)\), a set of ultrafilters of \(F(A)\), has at most
\(2^{\aleph_0}\leq\kappa\) points. Let \(P\) be the set of pairs
\((A,p)\) with \(A\subseteq B\) countable and \(p\in S(A)\); then
\(\kappa(P)\leq\kappa\). Add a new constant \(c_{(A,p)}\) for each pair
and let \(\Gamma\) be the diagram of \(B\) together with \(T\) and all
sentences \(\varphi(c_{(A,p)})\) with \(\varphi\in p\). A finite part of
\(\Gamma\) mentions finitely many new constants
\(c_1,\ldots,c_k\), for pairs \((A_i,p_i)\), and finitely many formulas of
each \(p_i\), whose conjunction \(\varphi_i\) lies in \(p_i\). Since
\(\varphi_i\) is consistent with the complete theory \(T(A_i)\), that
theory proves \(\exists v_0\,\varphi_i\). The model \(B\) satisfies
\(T(A_i)\), so some \(b_i\in B\) satisfies \(\varphi_i\); interpreting
\(c_i\) as \(b_i\) makes \(B\) a model of the finite part. By compactness
\(\Gamma\) has a model. Its language has at most \(\kappa\) symbols, so by
the downward Löwenheim--Skolem theorem it has a model of power at most
\(\kappa\); the interpretations of the names of \(B\) are distinct, so the
power is exactly \(\kappa\). The names give a monomorphism of \(B\) into
this model, because the diagram records every relation and its negation;
identifying \(B\) with its image yields \(B'\), and \(c_{(A,p)}\) realizes
\(p\).

\emph{The chain.} The theory has a model of power \(\kappa\) by the
Löwenheim--Skolem theorems. Let \(B_0\) be one, let \(B_{\alpha+1}\) be
obtained from \(B_\alpha\) by the step, and take unions at limits, for
\(\alpha<\omega_1\). Each \(B_\alpha\) is a model of \(T\) of power
\(\kappa\), since a union of a chain of models is a model and
\(\kappa\cdot\kappa(\alpha)=\kappa\). Let \(B=\bigcup_{\alpha<\omega_1}
B_\alpha\), a model of power \(\kappa\cdot\aleph_1=\kappa\), uncountable.
If \(A\subseteq B\) is countable, each element of \(A\) lies in some
\(B_\alpha\), and a countable set of countable ordinals has a bound below
\(\omega_1\); so \(A\subseteq B_\beta\) for some \(\beta<\omega_1\). Every
\(p\in S(A)\) is realized by some \(b\in B_{\beta+1}\). The inclusion
\(B_{\beta+1}\subseteq B\) is elementary, so \(b\) realizes \(p\) in
\(B\). Thus \(B\) is saturated over countable subsystems.
\end{proof}
\end{annexitem}

\begin{annexitem}{annex:p03-w03-p03-034-continuum-bound}{thm:p03-w03-p03-034-continuum-bound}{Embedding a countable system with continuum many types}
\begin{proof}
Since \(T\) is not totally transcendental, some countable
\(C\in\mathcal N(T)\) has \(\kappa(S(C))=2^{\aleph_0}\)
(\cref{thm:p03-morley-countability}); \(C\) is a subsystem of a model
\(M\). Let \(B\) be a model saturated over countable subsystems. We embed
\(C\) into \(B\). Enumerate \(|C|=\{c_n:n<\nu\}\) with \(\nu\leq\omega\),
let \(C_n\) be the induced system on \(\{c_m:m<n\}\), and build
increasing monomorphisms \(g_n:C_n\to B\), starting with the empty map.
Given \(g_n\), let \(p\) be the point of \(S(C_n)\) realized by \(c_n\) in
\(M\); its transport along \(g_n\) is a point over the finite subsystem
\(g_n(C_n)\) of \(B\), which is finite and hence countable, so it is
realized by some \(b\in B\), not in
\(g_n(C_n)\). By the transport criterion \(g_{n+1}=g_n\cup\{(c_n,b)\}\)
is a monomorphism. The union \(g\) is a monomorphism, and an isomorphism
of \(C\) onto the countable subsystem \(C'=g(C)\) of \(B\). The map
\(g^*\) is a bijection of \(S(C')\) onto \(S(C)\), so
\(\kappa(S(C'))=2^{\aleph_0}\). Every point of \(S(C')\) is realized in
\(B\), and distinct points by distinct elements, so
\(\kappa(B)\geq2^{\aleph_0}\).

If \(2^{\aleph_0}>\aleph_1\), then
\(\aleph_1^{\aleph_0}\geq2^{\aleph_0}>\aleph_1\), so \(\kappa=\aleph_1\)
satisfies \(\kappa^{\aleph_0}\neq\kappa\), and by the first part no model
of power \(\aleph_1\) of a theory that is not totally transcendental is
saturated over countable subsystems. \(\mathsf{DLO}\) is not totally
transcendental by \cref{thm:p03-morley-thin-models}, and saturation of
\(\Sigma\)-models and of their expansions coincide by
\cref{thm:p03-morley-normalization}.
\end{proof}
\end{annexitem}
